% GENERATED FILE. Edit canonical lessons, then run npm run generate:books.
% canonical-source-hash: fnv1a64:3cc3a19a856c21b1
% canonical-lessons: MR-C460-maya, MR-C460-jivhala, MR-C460-apulki, MR-C460-javlik, MR-C460-ladka

\chapter{Affection}
\label{ch:mr-affection}
\hlchaptermodality{460}

\begin{chapteropening}
\textbf{By the end of this chapter:} \emph{I can talk about affection, warmth and closeness.}

Complete the last lesson of chapter 460: I can talk about affection, warmth and closeness.
\end{chapteropening}

\section[\texorpdfstring{\mr{माया} (māyā)}{माया (māyā)}]{\mr{माया} (māyā) — affection; tender love}

\label{lesson:MR-C460-maya}



\begin{quote}
Before the new one: say the Marathi for to long for, then the Marathi for to pine.
\end{quote}

\subsection*{You'll want to know: \mr{माया}}
\textbf{\mr{माया}} — \emph{māyā} — \textquotedblleft{}affection; tender love\textquotedblright{}.

\begin{quote}
\mr{आजीची} \mr{माया} \mr{सगळ्या} \mr{मुलांवर} \mr{असते}.
\end{quote}

\begin{quote}
\emph{ājīcī māyā sagḷyā mulāṁvar aste.}
\end{quote}

\begin{quote}
Grandmother's affection is on all the children.
\end{quote}

Feminine.

Feminine; \mr{मायाळू}, loving. In other settings \mr{माया} also means illusion.

\begin{cousinweb}
\textbf{\mr{माया}} (\emph{māyā}) is a tatsama, taken from Sanskrit \emph{māyā} (CDIAL 10060), which CDIAL glosses \textquotedblleft{}supernatural power, skill; illusion\textquotedblright{} and, in the lexicons, \textquotedblleft{}compassion\textquotedblright{}.
\end{cousinweb}

\begin{grammarlens}[title={What you've built}]
One more word for describing what you lived and felt.
\end{grammarlens}

\subsection*{Your turn}
\begin{itemize}
\raggedright
  \item \emph{Say it:} \emph{māyā}
  \item \emph{Say it:} \emph{māyā}, once more
  \item \emph{Say it:} \emph{ājīcī māyā sagḷyā mulāṁvar aste.} — Grandmother's affection is on all the children.
  \item \emph{Recall:} say the Marathi for to long for, then the Marathi for to pine, then say \emph{māyā} again
\end{itemize}

\begin{tcolorbox}[breakable,colback=teal!4,colframe=teal!35!black,title={Before you move on}]
What does \mr{माया} mean? (\textquotedblleft{}Affection; tender love\textquotedblright{}.) Say it once more.
\end{tcolorbox}
\section[\texorpdfstring{\mr{जिव्हाळा} (jivhāḷā)}{जिव्हाळा (jivhāḷā)}]{\mr{जिव्हाळा} (jivhāḷā) — warmth of feeling}

\label{lesson:MR-C460-jivhala}



\begin{quote}
Before the new one: say the Marathi for to pine, then the Marathi for affection.
\end{quote}

\subsection*{You'll want to know: \mr{जिव्हाळा}}
\textbf{\mr{जिव्हाळा}} — \emph{jivhāḷā} — \textquotedblleft{}warmth of feeling\textquotedblright{}.

\begin{quote}
\mr{आमच्या} \mr{गल्लीतल्या} \mr{लोकांमध्ये} \mr{खूप} \mr{जिव्हाळा} \mr{आहे}.
\end{quote}

\begin{quote}
\emph{āmcyā gallītlyā lokāmmadhye khūp jivhāḷā āhe.}
\end{quote}

\begin{quote}
There is a lot of warmth among the people in our lane.
\end{quote}

Masculine.

Masculine; \mr{जिव्हाळ्याचा}, close to the heart: \mr{जिव्हाळ्याचा} \mr{मित्र}.

\begin{cousinweb}
\textbf{\mr{जिव्हाळा}} (\emph{jivhāḷā}) is a Marathi noun whose origin is not settled.
\end{cousinweb}

\begin{grammarlens}[title={What you've built}]
One more word for describing what you lived and felt.
\end{grammarlens}

\subsection*{Your turn}
\begin{itemize}
\raggedright
  \item \emph{Say it:} \emph{jivhāḷā}
  \item \emph{Say it:} \emph{jivhāḷā}, once more
  \item \emph{Say it:} \emph{āmcyā gallītlyā lokāmmadhye khūp jivhāḷā āhe.} — There is a lot of warmth among the people in our lane.
  \item \emph{Recall:} say the Marathi for to pine, then the Marathi for affection, then say \emph{jivhāḷā} again
\end{itemize}

\begin{tcolorbox}[breakable,colback=teal!4,colframe=teal!35!black,title={Before you move on}]
What does \mr{जिव्हाळा} mean? (\textquotedblleft{}Warmth of feeling\textquotedblright{}.) Say it once more.
\end{tcolorbox}
\section[\texorpdfstring{\mr{आपुलकी} (āpulkī)}{आपुलकी (āpulkī)}]{\mr{आपुलकी} (āpulkī) — fellow feeling}

\label{lesson:MR-C460-apulki}



\begin{quote}
Before the new one: say the Marathi for affection, then the Marathi for warmth of feeling.
\end{quote}

\subsection*{You'll want to know: \mr{आपुलकी}}
\textbf{\mr{आपुलकी}} — \emph{āpulkī} — \textquotedblleft{}fellow feeling\textquotedblright{}.

\begin{quote}
\mr{नव्या} \mr{शाळेत} \mr{सगळ्यांनी} \mr{माझं} \mr{आपुलकीने} \mr{स्वागत} \mr{केलं}.
\end{quote}

\begin{quote}
\emph{navyā śāḷet sagḷyānnī mājhaṁ āpulkīne svāgat kelaṁ.}
\end{quote}

\begin{quote}
At the new school everyone welcomed me warmly.
\end{quote}

Feminine.

Built on \mr{आपला}, one's own: treating someone as one's own.

\begin{cousinweb}
\textbf{\mr{आपुलकी}} (\emph{āpulkī}) is usually traced to \mr{आपुला}, an older form of \mr{आपला}, \textquotedblleft{}one's own\textquotedblright{}, with the noun ending \emph{-kī}.
\end{cousinweb}

\begin{grammarlens}[title={What you've built}]
One more word for describing what you lived and felt.
\end{grammarlens}

\subsection*{Your turn}
\begin{itemize}
\raggedright
  \item \emph{Say it:} \emph{āpulkī}
  \item \emph{Say it:} \emph{āpulkī}, once more
  \item \emph{Say it:} \emph{navyā śāḷet sagḷyānnī mājhaṁ āpulkīne svāgat kelaṁ.} — At the new school everyone welcomed me warmly.
  \item \emph{Recall:} say the Marathi for affection, then the Marathi for warmth of feeling, then say \emph{āpulkī} again
\end{itemize}

\begin{tcolorbox}[breakable,colback=teal!4,colframe=teal!35!black,title={Before you move on}]
What does \mr{आपुलकी} mean? (\textquotedblleft{}Fellow feeling\textquotedblright{}.) Say it once more.
\end{tcolorbox}
\section[\texorpdfstring{\mr{जवळीक} (javḷīk)}{जवळीक (javḷīk)}]{\mr{जवळीक} (javḷīk) — closeness}

\label{lesson:MR-C460-javlik}



\begin{quote}
Before the new one: say the Marathi for warmth of feeling, then the Marathi for fellow feeling.
\end{quote}

\subsection*{You'll want to know: \mr{जवळीक}}
\textbf{\mr{जवळीक}} — \emph{javḷīk} — \textquotedblleft{}closeness\textquotedblright{}.

\begin{quote}
\mr{एकत्र} \mr{प्रवास} \mr{केल्यावर} \mr{दोघींमध्ये} \mr{जवळीक} \mr{वाढली}.
\end{quote}

\begin{quote}
\emph{ekatra pravās kelyāvar doghīmmadhye javḷīk vāḍhlī.}
\end{quote}

\begin{quote}
After travelling together the two women grew close.
\end{quote}

Feminine.

Built on \mr{जवळ}, near: closeness between people.

\begin{cousinweb}
\textbf{\mr{जवळीक}} (\emph{javḷīk}) is usually traced to \mr{जवळ}, \textquotedblleft{}near\textquotedblright{}, with the noun ending \emph{-īk}.
\end{cousinweb}

\begin{grammarlens}[title={What you've built}]
One more word for describing what you lived and felt.
\end{grammarlens}

\subsection*{Your turn}
\begin{itemize}
\raggedright
  \item \emph{Say it:} \emph{javḷīk}
  \item \emph{Say it:} \emph{javḷīk}, once more
  \item \emph{Say it:} \emph{ekatra pravās kelyāvar doghīmmadhye javḷīk vāḍhlī.} — After travelling together the two women grew close.
  \item \emph{Recall:} say the Marathi for warmth of feeling, then the Marathi for fellow feeling, then say \emph{javḷīk} again
\end{itemize}

\begin{tcolorbox}[breakable,colback=teal!4,colframe=teal!35!black,title={Before you move on}]
What does \mr{जवळीक} mean? (\textquotedblleft{}Closeness\textquotedblright{}.) Say it once more.
\end{tcolorbox}
\section[\texorpdfstring{\mr{लाडका} (lāḍkā)}{लाडका (lāḍkā)}]{\mr{लाडका} (lāḍkā) — beloved; pet}

\label{lesson:MR-C460-ladka}



\begin{quote}
Before the new one: say the Marathi for fellow feeling, then the Marathi for closeness.
\end{quote}

\subsection*{You'll want to know: \mr{लाडका}}
\textbf{\mr{लाडका}} — \emph{lāḍkā} — \textquotedblleft{}beloved; pet\textquotedblright{}.

\begin{quote}
\mr{ही} \mr{मांजर} \mr{आमची} \mr{लाडकी} \mr{आहे}.
\end{quote}

\begin{quote}
\emph{hī māñjar āmcī lāḍkī āhe.}
\end{quote}

\begin{quote}
This cat is our darling.
\end{quote}

Agrees: \mr{लाडका} \mr{मुलगा}, \mr{लाडकी} \mr{मुलगी}, \mr{लाडकं} \mr{बाळ}.

\begin{cousinweb}
\textbf{\mr{लाडका}} (\emph{lāḍkā}) is inherited: CDIAL lists Marathi \emph{lāḍkā}, \textquotedblleft{}darling\textquotedblright{}, under the reconstructed \emph{*laḍikka}, \textquotedblleft{}child\textquotedblright{} (CDIAL 10924).
\end{cousinweb}

\begin{grammarlens}[title={What you've built}]
One more word for describing what you lived and felt.
\end{grammarlens}

\subsection*{Your turn}
\begin{itemize}
\raggedright
  \item \emph{Say it:} \emph{lāḍkā}
  \item \emph{Say it:} \emph{lāḍkā}, once more
  \item \emph{Say it:} \emph{hī māñjar āmcī lāḍkī āhe.} — This cat is our darling.
  \item \emph{Recall:} say the Marathi for fellow feeling, then the Marathi for closeness, then say \emph{lāḍkā} again
\end{itemize}

\begin{tcolorbox}[breakable,colback=teal!4,colframe=teal!35!black,title={Before you move on}]
What does \mr{लाडका} mean? (\textquotedblleft{}Beloved; pet\textquotedblright{}.) Say it once more.
\end{tcolorbox}
